\section{Valoración de la corrección del diseño}

El diseño es correcto en su estructura y responde al problema. Cubre los actores, separa la autoría del bloque de la generación del horario, guarda los parámetros junto al horario, modela la validación del coordinador, la aprobación del jefe, la asignación, la asistencia y el cambio con coordinador de apoyo, y saca el histórico del camino transaccional. Esa separación es la decisión de diseño que el enunciado exige de forma explícita.

\subsection{Lo que el modelo representa bien}

\begin{itemize}
  \item Cada requerimiento informacional de la sección~3 tiene entidad, atributo o relación. Los seis reportes del enunciado se pueden armar con horario, creador, parámetros, validación, bloque, complejidad, asignatura, asistencia y grupo.
  \item La especialización deja claro que administrador, coordinador, estudiante y jefe no son el mismo papel.
  \item CREA, GENERA, la validación y la aprobación no se colapsan en un solo «responsable del horario». La consulta de revisores lee VALIDACIÓN. El horario final sale de la decisión aprobado, que escribe el estado.
  \item La sesión ubica el bloque en una franja del catálogo y en un aula. La fecha de cada impartición está en la entidad ASISTENCIA, porque el hueco semanal se repite. La asistencia no duplica el horario.
  \item Las relaciones N:N reales (autorización, matrícula y asignación) no se meten como listas dentro de una columna. La asignatura no es N:N: pertenece a un programa. El grupo tampoco es N:N ni entidad: es un atributo del estudiante.
  \item Los atributos derivados (hash_duplicado, equilibrio_cumplido, tasas del analítico) están marcados como derivados. Guardarlos es una decisión de consulta, no una confusión con un dato fuente. El duplicado lo cierra la unicidad de las columnas fuente. El equilibrio lo recalcula un disparador.
  \item El OLTP de negocio, quitando esos derivados y la proyección de lectura, está en tercera forma normal. La proyección no es una relación de negocio: la reconstruyen los disparadores desde las sesiones.
  \item El esquema estrella responde sumas, promedios y comparaciones por periodo, aula y día. El hecho de asistencia conserva estudiante, grupo, programa y el bloque con su complejidad y su creador. El hecho de solicitud, en el mismo almacén, identifica a quienes pidieron cambio. El promedio del grupo y los diez bloques con más inasistencia salen de ahí, sin volver al operacional.
\end{itemize}

\subsection{Conclusión}

El MERX operacional, leído con la participación de la sección~\ref{sec:relaciones}, y el esquema estrella separado son un diseño correcto para el sistema de horarios docentes: completo respecto a los requerimientos reformulados, sin mezclar la operación con el histórico, y normalizado en la parte transaccional.

El DDL operacional y analítico cierra el solape de sesiones, el choque de aula por intervalo, el catálogo de franjas sin cruces, el bloque del mismo programa, el duplicado de bloque, el equilibrio recalculado, la matrícula previa a la asignación, la validación del coordinador, una sola aprobación del jefe con sucesor del mismo programa, la coherencia de la solicitud y la asignación que produce un cambio aceptado con destino. No cierra la asignación con una unicidad por estudiante, programa y semestre: el enunciado no la pide. Quedan en la aplicación la clasificación exclusiva del creador del bloque, el filtro de horarios aprobados en los reportes, la copia de los valores de parámetros al horario sucesor y la exportación a PDF con el orden de columnas.
